\section{Introduction}

Phase 1 of this project established a controlled reference workflow for spike-based DOA-sector inference in an embedded-AI context. The completed prototype used a physics-based ULA signal model, sample covariance matrix estimation, beamspace feature extraction, nearest-centroid baseline classification, rate-based spike encoding, and LIF-inspired temporal inference. The outcome was not merely a Python classification script; it was a staged engineering baseline connecting radar/array signal processing to sparse temporal inference.

The purpose of Phase 2 is to extend this baseline without breaking continuity. The new phase will ask whether the Phase 1 pipeline remains meaningful when the radar/DOA sensing conditions become more challenging. Instead of immediately moving to FPGA deployment or fixed-point hardware synthesis, Phase 2 will strengthen the scientific foundation of the sensing and inference problem. This is necessary because hardware acceleration is only meaningful after the algorithmic behavior is understood under realistic input conditions.

The proposed Phase 2 title is:

\begin{proposalbox}{Proposed Phase 2 Title}
\textbf{Robust Radar/DOA Sensing Case Study Based on the Phase 1 Spike-Inference Pipeline}
\end{proposalbox}

This title makes the continuity explicit. Phase 2 is not a separate project. It is a robustness and extension study built directly on the Phase 1 result.
